% Part: turing-machines
% Chapter: undecidability
% Section: universal-tm

\documentclass[../../../include/open-logic-section]{subfiles}

\begin{document}

\olfileid{tur}{und}{uni}
\olsection{نړيوال ټيورينګ ماشينونه}

په~\olref[enu]{sec} کښې مو وڅېړل چې هر ټيورينګ ماشين د صحيح اعدادو
په يوه متناهي لړۍ بيانېدے شي۔ دا لړۍ د ټيورينګ ماشين حالتونه، الفبا،
پيلنی حالت، او لارښوونې کوډوي۔ موږ دا هم په ګوته کړه چې د دې ټولو
بيانونو سټ !!{enumerable} دے۔ څرنګه چې د داسې بيانونو سټ
!!{denumerable} دے، معنا ئې دا ده چې له~$\Nat$ څخه دې بيانونو ته يوه
!!a{surjective} تابع شته۔ داسې !!a{surjective} تابع مثلاً د کانتور د
زېګ‌زېګ ميتود په کارولو تر لاسه کېدے شي۔ دا موږ ته د ټيورينګ
ماشينونو د ټولو (بيانونو) د شمېرلو لار راکوي۔ که يوه داسې شمېره
ثابته کړو، نو اوس د~$1$-م، $2$-م، \dots، او~$e$-م ټيورينګ ماشين په
اړه خبرې معنا لري۔ دې عددونو ته \emph{شاخصونه} ويل کېږي۔

\begin{defn}
که~$M$ (زموږ په ثابته شوې شمېره کښې) $e$-م ټيورينګ ماشين وي، وايو
چې~$e$ د~$M$ يو \emph{شاخص} دے۔ د~$e$-م ټيورينګ ماشين دپاره~$M_e$
ليکو۔
\end{defn}

يو ماشين ښايي له يوه څخه ډېر شاخصونه ولري؛ مثلاً د~$M$ دوه بيانونه
ښايي د هغو د لارښوونو د لست کولو په ترتيب کښې توپير ولري، او دا
بېلابېل بيانونه به بېلابېل شاخصونه ولري۔

مهمه دا ده چې د ټيورينګ ماشين د بيانونو شمېره داسې ورکول کېدے شي
چې موږ د~$M$ بيان د هغه له شاخص څخه، او د يوه ماشين~$M$ يو شاخص د
هغه له بيان څخه په اغېزمنه توګه محاسبه کړے شو۔ بيا د چرچ--ټيورينګ د
اصل له مخې داسې ټيورينګ ماشين موندل کېدے شي چې د شاخص~$e$ لرونکي
ټيورينګ ماشين بيان بېرته تر لاسه کړي او اړوند بيان پر خپله پټه د وت
په توګه وليکي۔ بيان به د~$\TMstroke$ د بلاکونو يوه لړۍ وي (چې د~$M_e$
په بيانوونکې لړۍ کښې مثبت صحيح اعداد تمثيلوي)۔

دې ته په پام اوس دا پوښتنه طبيعي ده: د ټيورينګ ماشين د شاخصونو کومې
تابعې پخپله د ټيورينګ ماشينونو په وسيله محاسبه کېدونکې دي؟ د ټيورينګ
ماشين د شاخصونو کوم خاصيتونه د ټيورينګ ماشينونو په وسيله پرېکړه کېدے
شي؟ يوه بېلګه: هغه تابع چې شاخص~$e$ د هغه شاخص~$e$ لرونکي ټيورينګ ماشين
د حالتونو شمېر ته وړي، د يوه ټيورينګ ماشين په وسيله محاسبه کېدونکې
ده۔ داسې ټيورينګ ماشين به دا کار کوي: پر هغې پټې له پيل وروسته چې د
$e$ دانو~$\TMstroke$ يو بلاک لري، لومړے به~$e$ په خپل بيان ناکوډ
کړي۔ بيان اوس پر پټه د~$\TMstroke$ د بلاکونو په يوه لړۍ تمثيلېږي۔
څرنګه چې په دې لړۍ کښې لومړے !!{element} د حالتونو شمېر دے، اوس
يوازې دا پاتې دي چې د~$\TMstroke$ له لومړي بلاک پرته نور هر څه پاک
کړي او بيا ودرېږي۔

يوه د پام وړ پايله دا ده:

\begin{thm}\ollabel{thm:universal-tm} داسې يو \emph{نړيوال ټيورينګ
  ماشين}~$U$ شته چې کله پر ننوت~$\tuple{e,n}$ پيل شي،
  \begin{enumerate}
    \item هله او يوازې هله درېږي چې~$M_e$ پر ننوت~$n$ ودرېږي، او
    \item که~$M_e$ له وت~$m$ سره ودرېږي، نو~$U$ هم همداسې کوي۔
  \end{enumerate}
  نو~$U$ هغه تابع~$f\colon \Nat \times \Nat \pto \Nat$ محاسبه کوي
  چې~$f(e,n) = m$ وي که~$M_e$ پر ننوت~$n$ له پيل وروسته له وت~$m$
  سره ودرېږي، او که نۀ وي نو ناتعريفه ده۔
\end{thm}

\begin{proof}
  په عمل کښې د~$U$ جوړول بنسټيز ډول ناشوني دي، ځکه دا يو بېخي پېچلے
  ماشين دے۔ خو موږ د هغه د کار عمومي طرحه بيانولے، او بيا د
  چرچ--ټيورينګ اصل رابللے شو۔ کله چې~$U$ پيل شي، پټه ئې د~$e$ دانو
  $\TMstroke$ يو بلاک لري چې ورپسې د~$n$ دانو~$\TMstroke$ يو بلاک
  دے۔ لومړے شاخص~$e$ د ننوت~$n$ ښي لور ته ``ناکوډوي''۔ له دې څخه د
  اعدادو يو لست جوړېږي (يعنې د~$\TMstroke$ هغه بلاکونه چې په
  $\TMblank$ سره بېلېږي) چې د~$M_e$ ماشين لارښوونې بيانوي۔ بيا~$U$
  د~$M_e$ د پيلني حالت عدد او عدد~$1$ د~$M_e$ د بيان ښي لور ته پر
  پټه ليکي۔ (دا هم په يوګوني ډول، د~$\TMstroke$ د بلاکونو په توګه،
  تمثيلېږي۔) ورپسې ننوت (د~$n$ دانو~$\TMstroke$ بلاک) ښي لور ته
  کاپي کوي---خو هر~$\TMstroke$ د درې دانو~$\TMstroke$ په يوه بلاک
  بدلوي (په ياد ولرئ، د~$\TMstroke$ نښې عدد~$3$ دے، $1$ د~$\TMendtape$
  عدد او~$2$ د~$\TMblank$ عدد دے)۔ د دې بلاکونو د لړۍ په کيڼ سر کښې
  (چې د~$U$ پر پټه د~$\TMblank$ نښې سره بېلېږي) يو~$\TMstroke$، يعنې
  د~$\TMendtape$ کوډ، ليکي۔

  اوس د~$U$ پر پټه دا شيان دي: شاخص~$e$، عدد~$n$، د پيلني حالت کوډ
  عدد (``اوسنی حالت'')، د سر د پيلني موقعيت عدد~$1$ (``د سر اوسنی
  موقعيت'')، او د ``پټې'' پيلنۍ منځپانګه (د~$\TMstroke$ د بلاکونو
  يوه لړۍ چې د~$M_e$ د نښو کوډ عددونه---``نښې''---تمثيلوي او په
  $\TMblank$ سره بېلېږي)۔

  اوس هغه څه په دې ډول مشابه چلوي چې~$M_e$ به پر ننوت~$n$ له پيل
  وروسته کړي وے:
  \begin{enumerate}
    \item د ``سر د اوسني موقعيت'' عدد~$k$ پيدا کول (په پيل کښې دا~$1$
    دے)،
    \item د ``پټې'' $k$-م بلاک ته تلل، څو وګوري چې هلته کومه ``نښه''
    ده،
    \item\ollabel{find-inst}%
    له اوسني ``حالت'' او ``نښې'' سره سمه لارښوونه پيدا کول،
    \item د ``پټې'' $k$-م بلاک ته بېرته تلل او هلته ``نښه'' د هغې
    نښې په کوډ عدد بدلول چې~$M_e$ به ليکلې وے،
    \item سر هغه ځاے ته وړل چې اوسنی ``حالت'' پکې ثبتېږي، او هلته
    عدد د نوي حالت په عدد بدلول،
    \item هغه ځاے ته تلل چې د ``پټې موقعيت'' پکې ثبتېږي، او يو
    $\TMstroke$ پاکول يا يو~$\TMstroke$ زياتول (په ترتيب سره که
    لارښوونه د کيڼ يا ښي لور د حرکت خبره کوي)۔
    \item تکرارول۔\footnote{موږ دلته پر ځينو نازکو ستونزو سرسري
    تېرېږو۔ مثلاً کله چې~$U$ هغه شمېرند زياتوي چې د ``سر اوسنی
    موقعيت'' پرې ساتي، ښايي څه زيات ځاے ته اړ شي---په هغه حالت کښې
    به ټوله ``پټه'' ښي لور ته وخوځوي۔}
  \end{enumerate}
  که~$M_e$ پر ننوت~$n$ له پيل وروسته هېڅکله ونۀ درېږي، نو~$U$ هم
  هېڅکله نۀ درېږي، ځکه وت ئې ناتعريفه دے۔

  که په ګام~\olref{find-inst} کښې څرګنده شي چې د~$M_e$ بيان د اوسني
  ``حالت''/``نښې'' جوړې دپاره هېڅ لارښوونه نۀ لري، نو~$M_e$ به
  درېدلے وے۔ که دا پېښه شي، $U$ د خپلې پټې هغه برخه پاکوي چې د
  ``پټې'' کيڼ لور ته ده۔ د درې دانو~$\TMstroke$ د هر بلاک دپاره
  (چې د~$M_e$ پر پټه يو~$\TMstroke$ تمثيلوي) د خپلې پټې پر کيڼ سر
  يو~$\TMstroke$ ليکي، او ``پټه'' پرله‌پسې پاکوي۔ کله چې دا کار بشپړ
  شي، د~$U$ پټه د~$m$ اوږدوالي د~$\TMstroke$ يوازې يو بلاک لري۔

  که~$U$ پر ``پټه'' د درې دانو~$\TMstroke$ له بلاک پرته له بل څه
  سره مخ شي، سملاسي درېږي۔ څرنګه چې په دې حالت کښې د~$U$ پټه د
  $\TMstroke$ يوازې يو بلاک نۀ لري، وت ئې طبيعي عدد نۀ دے؛ يعنې په
  دې حالت کښې~$f(e,n)$ ناتعريفه ده۔
\end{proof}

\end{document}
